\documentclass[12pt]{article}
\usepackage[utf8]{inputenc}
\usepackage[T1]{fontenc}
\usepackage{hyperref}
\title{Teke-Tege Language Resources for K-12 Curriculum Development}
\author{Compiled by Claude Sonnet 5 (claude-sonnet-5) for Andr\'as Kornai}
\date{}
\begin{document}
\maketitle

\section{Introduction}
Teke-Tege (ISO 639-3 teg; also called North Teke, Latege or Tege) is a Bantu language of southeastern Gabon and eastern Congo, with about 140,000 speakers. Do not confuse it with Tsege (tck). The searches found a small dictionary and some grammar studies, but no school materials.

\section{Primary Educational Resources}
No primers, textbooks or school programmes were found.

\section{Linguistic and Reference Materials}
\begin{itemize}
\item Latege dictionary by Pauline Linton on Webonary. The site says it is at an early stage: \url{https://www.webonary.work/latege/overview/introduction/?lang=en}
\item Linton, \emph{Phrases non verbales du latege} (PDF, in French): \url{https://www.webonary.work/latege/files/Phrases-non-verbales-du-latege.pdf}
\item Linton, \emph{Classes nominales et syntagmes nominaux du latege} (PDF, in French): \url{https://www.webonary.work/latege/files/Classes-nominales-et-syntagmes-nominaux-du-latege.pdf}
\item North Teke language, Wikipedia overview: \url{https://en.wikipedia.org/wiki/North_Teke_language}
\item Omniglot page on North Teke: \url{https://omniglot.com/writing/northteke.htm}
\item Glosbe Teke-Tege entries. This is user-contributed, so treat its quality with caution: \url{https://glosbe.com/teg/lib}
\item ASJP wordlist for Northern Teke: \url{https://asjp.clld.org/languages/NORTHERN_TEKE}
\end{itemize}

\section{Religious and Community Resources}
None were found.

\section{Digital and Audio Resources}
The only digital resource found is the Webonary dictionary above. No audio was found.

\section{Curriculum Development Context}
The language is not known to be taught in schools. It is written in the Latin script, and the Linton orthography is used in the Webonary dictionary. Building a curriculum would depend on extending the Linton materials with community authors.

\end{document}